% Experiments Section Update - PDB Validation Results
% Add to Results section after existing validation content

\subsection{PDB Structure Comparison Validation}
\label{sec:pdb_validation}

To provide independent validation of stereochemistry correctness, we performed systematic comparison with experimentally determined PDB structures containing verified glycan stereochemistry.

\subsubsection{Validation Method}

We selected four PDB structures representing diverse glycan types with critical stereochemistry features:

\begin{table}[h]
\centering
\caption{PDB Reference Structures for Validation}
\begin{tabular}{lllp{4cm}}
\toprule
ID & PDB & Structure & Critical Validation Point \\
\midrule
V001 & 5NSC & Fucosylated & FUC = alpha-L-fucose (NOT beta) \\
V002 & 1L6R & Lactose & GAL C4 axial hydroxyl (epimer) \\
V003 & 2VXR & Sialyllactose & SIA anomeric = C2 (ketose) \\
V004 & 5K65 & M3 N-glycan & Branch topology preservation \\
\bottomrule
\end{tabular}
\end{table}

For each reference structure, we:
\begin{enumerate}
\item Extracted glycan IUPAC notation from PDB annotations
\item Processed through GlycoSMILES2BAP CCD mapper
\item Compared output CCD codes with PDB-deposited CCD codes
\item Verified anomeric carbon positions and ring oxygen assignments
\end{enumerate}

\subsubsection{Validation Results}

All four test cases passed validation (100\% success rate):

\begin{table}[h]
\centering
\caption{PDB Validation Results}
\begin{tabular}{llccccl}
\toprule
Test & PDB & Expected & Actual & Anomeric & Ring O & Status \\
\midrule
Fucosylated & 5NSC & FUC & FUC & C1 & O5 & \textbf{PASS} \\
Lactose & 1L6R & GAL & GAL & C1 & O5 & \textbf{PASS} \\
Sialyllactose & 2VXR & SIA & SIA & \textbf{C2} & \textbf{O6} & \textbf{PASS} \\
M3 N-glycan & 5K65 & MAN/BMA/NAG & MAN/BMA/NAG & C1 & O5 & \textbf{PASS} \\
\bottomrule
\end{tabular}
\end{table}

\textbf{Critical Finding - Sialic Acid C2 Anomeric:} The most important validation confirmed correct handling of Neu5Ac (sialic acid):
\begin{itemize}
\item Anomeric carbon correctly identified as C2 (not C1)
\item Ring oxygen correctly identified as O6 (not O5)
\item This addresses the core stereochemistry problem identified by Huang et al.\ (2025)
\end{itemize}

\textbf{Fucose L-Configuration:} The tool correctly distinguished:
\begin{itemize}
\item FUC (alpha-L-fucose) for the 5NSC structure
\item NOT BDF (beta-L-fucose) - wrong anomer
\item L-configuration preserved through mapping
\end{itemize}

\textbf{Epimer Distinction:} For lactose (1L6R):
\begin{itemize}
\item Galactose correctly mapped to GAL
\item NOT GLC (glucose) - different C4 stereochemistry
\item Epimer distinction validated
\end{itemize}

\textbf{Branch Topology:} For M3 N-glycan (5K65):
\begin{itemize}
\item Branch topology correctly preserved
\item All five residues correctly assigned
\item Alpha/beta distinction maintained (MAN vs BMA)
\end{itemize}

\subsubsection{Comparison with Literature-Reported Errors}

The PDB validation directly addresses documented stereochemistry errors:

\begin{table}[h]
\centering
\caption{Literature Error Comparison}
\begin{tabular}{llll}
\toprule
Source & Problem & Our Tool & Status \\
\midrule
Huang et al.\ 2025 & SMILES $\sim$60\% error rate & BAP format & Corrected \\
Frenz et al.\ 2018 & PDB:5NSC wrong anomer & FUC mapped & Corrected \\
Jo et al